% 张三慧《学习辅导与习题解答》扫描 S 卷1第11章“思考题选答”原编号。
% 仅辅导书来源已视觉核验，V1主教材原题逐图交叉核验仍待；全部exam_ready=false。
\section*{第11章\quad 热力学第二定律：辅导书已列思考题独立解析}
\paragraph{思考题 11.4}
固定物质的量理想气体等温线为 \(pV=nRT_0\)，可逆绝热线满足 \(pV^\gamma=C\) 且 \(\gamma>1\)。消去p得 \(nRT_0 V^{\gamma-1}=C\)，当 \(V>0\) 时只对应一个正V解，故\(\boxed{\text{两线最多相交一次}}\)。熵法也可说明理想气体等温过程 \(dS=nR\,dV/V\) 严格单调，绝热线熵不变。对具有异常热膨胀、相变等复杂状态方程的材料不可无条件套用“任何等温线与绝热线最多一次”这一未经限定的断言。

\paragraph{思考题 11.6}
熵增加原理要求\emph{整个孤立系统} \(\Delta S_{\rm total}\ge0\)，而水本身与空气有热交换，并不孤立。若水在 \(T_h\) 附近向环境 \(T_c<T_h\) 传出微量热 \(Q>0\)，有 \(\Delta S_{\rm water}\simeq-Q/T_h<0\)，空气熵变 \(\Delta S_{\rm air}\simeq Q/T_c>0\)，总变 \(\boxed{\Delta S_{\rm total}\simeq Q(1/T_c-1/T_h)>0}\)。有有限温差传热就存在正熵产生，局部降熵与第二定律不矛盾。

\paragraph{思考题 11.8}
取能级 \(E_{\rm low}<E_{\rm high}\)、简并度相同的有限自旋系统，平衡占据比 \(N_{\rm high}/N_{\rm low}=\exp[-(E_{\rm high}-E_{\rm low})/(k_BT_s)]\)。若高能态发生粒子数反转 \(N_{\rm high}>N_{\rm low}\)，指数只能由 \(\boxed{T_s<0}\) 描述。因自旋能谱有上界，熵 \(S(E)\) 在中点以上随能量上升反而下降，使 \(\partial S/\partial E=1/T_s<0\)。这样的负温自旋子系统与任意正温系统接触会趋向放出能量，故“热度”高于所有正温系统。真实晶格平动、振动能谱无上界，不能把整个晶体普通热平衡温度直接叫负温。
